\begin{abstract}
Personal digital twins built from conversational histories are often described in terms of persona, memory, or stylistic imitation, but these descriptions do not directly test whether a system can predict a user's future conversational behavior. We introduce a proxy replay benchmark for personal GPT chat logs: given held-out contexts from a user's historical conversation, predict the user's next conversational action under leakage-free evaluation. After tightening a rule-based topic-shift label, a lexical nearest-neighbor baseline using CJK-aware character n-gram similarity achieves the highest point estimates on a fixed chronological split (0.35 accuracy, 0.38 macro-F1), while the learned state router scores 0.23 accuracy and 0.30 macro-F1. We do not detect a significant paired accuracy difference (McNemar $p=0.074$), but bootstrap 95\% CIs are wide (lexical macro-F1 [0.15, 0.53]) given $n=40$ examples. A direct ablation shows that enabling state routing on top of the lexical baseline contributes zero additional signal — 40/40 predictions identical; this null result replicates across two pretrained embedding models (BGE-small-zh, MiniLM). Neither embedding model exceeds the CJK lexical tokenizer. The state-prior baseline confirms that dominant-action priors inflate accuracy while collapsing macro-F1. These results show that explicit state routing can appear useful under weak baselines but contributes nothing once lexical similarity is properly measured. The benchmark is single-user, rule-labeled, and not evidence of a complete human digital twin.
\end{abstract}
